% =============================================================
% Capítulo — Desenvolvimento
% =============================================================
% Reescrito em 22/08/2026: o conteúdo anterior descrevia o experimento
% com ResNet-50 / CIFAR-100, anterior ao pivô para Transformer decidido
% na reunião de 06/05/2026. Esta versão descreve o experimento
% efetivamente implementado, com GPT-2 small sobre o WikiText-2.
% =============================================================

\chapter{Desenvolvimento}\label{chap:desenvolvimento}

O Capítulo~\ref{chap:metodologia} apresentou os critérios que orientam o percurso
metodológico deste trabalho. Este capítulo descreve as decisões concretas que decorrem
desses critérios e o experimento tal como foi efetivamente implementado: o corpus e o
modelo selecionados, quais parâmetros são considerados podáveis, como cada estratégia de
poda foi implementada, como o perfil de sensibilidade é levantado e convertido em taxas
por camada, o protocolo de \textit{fine-tuning} de recuperação, o \textit{framework} e as
bibliotecas empregadas, a organização do \textit{pipeline} experimental, as métricas
coletadas, o ambiente de execução e as limitações reconhecidas.

\section{Visão Geral}\label{sec:visao_geral}

A pergunta que orienta o trabalho é empírica e comparativa. Não se pretende demonstrar que
a poda comprime modelos --- resultado já consolidado pela literatura
\cite{Han2015, Li2017, Blalock2020} ---, mas responder a duas questões mais específicas,
sob condições controladas e reprodutíveis:

\begin{enumerate}
    \item[(A)] Em que medida a redução \textit{teórica} de parâmetros se converte em
          redução \textit{real} de tempo de inferência e de consumo de energia, e como
          essa conversão depende da granularidade da poda?
    \item[(B)] Uma alocação não uniforme do orçamento de poda, guiada por uma análise de
          sensibilidade por camada, melhora o compromisso entre compressão e qualidade em
          relação às alocações ingênuas?
\end{enumerate}

Responder a essas duas perguntas exige que modelos comprimidos por estratégias distintas
sejam avaliados exatamente pelo mesmo conjunto de métricas, no mesmo corpus e no mesmo
\textit{hardware}. Essa simetria de avaliação é condição necessária para que qualquer
comparação seja interpretável, e é o princípio que organiza todo o desenvolvimento
descrito adiante.

\subsection{As Quatro Etapas Centrais}\label{subsec:quatro_etapas}

O experimento é organizado em quatro etapas sequenciais, cada uma com papel distinto.

A \textbf{primeira etapa} estabelece o modelo de referência. O GPT-2 \textit{small}
pré-treinado é carregado sem qualquer modificação e avaliado no conjunto de teste do
WikiText-2, registrando-se o vetor completo de métricas --- perplexity, contagem de
parâmetros, tamanho em memória, FLOPs, tempo de inferência e consumo de energia. Esse
modelo denso é o referencial absoluto contra o qual todas as configurações comprimidas
serão comparadas.

A \textbf{segunda etapa} aplica as duas estratégias clássicas de poda, de forma
independente e sempre a partir do modelo denso original: a poda não estruturada por
magnitude \cite{Han2015} e a poda estruturada por norma $L_1$ \cite{Li2017}, esta última
adaptada ao \textit{Transformer}. Para cada estratégia, varre-se toda a faixa de
esparsidade, de 10\% a 90\% em passos de 10 pontos percentuais, e sob dois escopos de
alocação distintos. A escolha dessas duas estratégias é deliberada: elas ocupam os dois
extremos do espectro de granularidade e, como se verá, comportam-se de forma oposta quanto
ao impacto real sobre o \textit{hardware}.

A \textbf{terceira etapa} introduz a estratégia otimizada proposta, que constitui a
contribuição central do trabalho: a análise de sensibilidade por camada. Ela se divide em
duas fases. Na primeira, levanta-se o perfil de sensibilidade do modelo, podando uma
camada de cada vez e medindo a degradação resultante. Na segunda, esse perfil é convertido
em uma alocação não uniforme de taxas por camada, que é então avaliada nos mesmos níveis
de orçamento das varreduras ingênuas.

A \textbf{quarta etapa} aplica o \textit{fine-tuning} de recuperação. Todas as etapas
anteriores operam em regime \textit{one-shot}, isto é, podam e medem sem retreinamento, de
modo a isolar o dano causado pela poda. Esta etapa reintroduz o retreinamento --- como
fazem tanto \citeonline{Han2015} quanto \citeonline{Li2017} --- e mede quanto da
degradação é recuperável.

\subsection{Discussão das Decisões Metodológicas}\label{subsec:discussao_decisoes}

Todo desenho experimental implica escolhas que condicionam o escopo dos resultados. A
Tabela~\ref{tab:decisoes_metodologicas} sintetiza as principais decisões adotadas, as
alternativas consideradas e a justificativa de cada uma. As seções seguintes desenvolvem
em maior profundidade as que são mais críticas para a interpretação dos resultados.

\begin{table}[htb]
\centering
\caption{Síntese das principais decisões metodológicas do trabalho.}
\label{tab:decisoes_metodologicas}
\small
\begin{tabular}{p{2.6cm} p{3.4cm} p{2.6cm} p{4.4cm}}
\hline
\textbf{Decisão} & \textbf{Alternativas consideradas} & \textbf{Escolha} & \textbf{Justificativa} \\
\hline

\textit{Dataset}
& WikiText-103, The Pile, OpenWebText
& \textbf{WikiText-2}
& Referência consolidada para perplexity de modelos de linguagem, com custo de avaliação
  compatível com o \textit{hardware} disponível. \\

Modelo
& GPT-2 \textit{medium}/\textit{large}, BERT, modelos da família LLaMA
& \textbf{GPT-2 \textit{small}}
& Arquitetura \textit{Transformer} representativa, com pesos públicos e ampla presença na
  literatura de compressão; 124M parâmetros cabem na GPU disponível. \\

Estratégias de poda
& Três ou mais estratégias
& \textbf{Duas} (magnitude e estruturada)
& Cobrem os dois extremos de granularidade; restrição de tempo experimental
  inviabilizaria uma terceira com a mesma profundidade. \\

Regime de poda
& Iterativo (podar e retreinar em ciclos)
& \textbf{\textit{One-shot} primeiro}
& Isola o efeito da poda do efeito do retreinamento; o \textit{fine-tuning} entra depois,
  como etapa separada e mensurável. \\

Estratégia otimizada
& Destilação de conhecimento, quantização
& \textbf{Sensibilidade por camada}
& Ataca diretamente a hipótese do trabalho --- a de que as camadas não são igualmente
  importantes --- sem alterar a natureza da poda. \\

Métrica principal
& \textit{Accuracy} em tarefas descendentes
& \textbf{Perplexity}
& Métrica intrínseca padrão para modelos de linguagem, mensurável sem
  \textit{fine-tuning} em tarefa específica. \\

Ambiente
& Máquina local, Google Colab
& \textbf{Kaggle (GPU T4)}
& Acesso gratuito a GPU dedicada com sessões longas o bastante para as varreduras
  completas. \\
\hline
\end{tabular}\\
\textit{Fonte: Elaboração própria.}
\end{table}

\section{\textit{Dataset}}\label{sec:dataset}

\subsection{O WikiText-2}\label{subsec:wikitext2}

O WikiText-2 é um corpus de língua inglesa extraído de artigos verificados da Wikipédia,
amplamente utilizado como referência na avaliação de modelos de linguagem. Diferentemente
de corpora construídos a partir de frases isoladas, o WikiText preserva a estrutura
original dos artigos --- parágrafos completos, pontuação e capitalização mantidas ---, o
que o torna adequado para medir a capacidade de um modelo de acompanhar contexto ao longo
de textos extensos.

Emprega-se especificamente a variante \texttt{wikitext-2-raw-v1}, na qual o texto não
sofre substituição de palavras raras por um marcador genérico. Essa escolha é necessária
porque o GPT-2 utiliza tokenização por \textit{subwords}, capaz de representar palavras
fora do vocabulário decompondo-as em fragmentos; a variante processada anularia essa
capacidade e distorceria a perplexity medida.

\subsection{Justificativa da Escolha}\label{subsec:dataset_justificativa}

Três fatores motivaram a escolha. O primeiro é a comparabilidade: perplexities medidas no
WikiText-2 são reportadas em grande número de trabalhos, o que permite verificar se o
\textit{baseline} obtido é coerente com o esperado --- verificação efetivamente realizada,
como se discute na Seção~\ref{sec:res_protocolo}. O segundo é o custo: o conjunto de teste
é pequeno o suficiente para que uma avaliação completa de perplexity leve menos de um
minuto na GPU utilizada, o que é determinante em um trabalho que executa mais de uma
centena de avaliações. O terceiro é a existência de uma divisão oficial em treino,
validação e teste, que dispensa decisões arbitrárias de particionamento.

Registre-se que o corpus é parametrizável na implementação: a substituição pelo
WikiText-103, de mesma natureza e cerca de cinquenta vezes maior, exige apenas a alteração
de um argumento de linha de comando. Essa flexibilidade foi mantida deliberadamente para
que trabalhos futuros possam verificar a generalidade dos resultados sem reescrever o
código.

\subsection{Pré-processamento e Tokenização}\label{subsec:preprocessamento}

O pré-processamento é intencionalmente mínimo, para que a perplexity medida reflita o
corpus e não decisões de limpeza. As linhas do conjunto são unidas por quebra dupla,
preservando a separação entre parágrafos tal como no corpus original, e o texto resultante
é tokenizado como um fluxo contínuo único.

A tokenização emprega o tokenizador original do GPT-2, baseado no algoritmo BPE
(\textit{Byte Pair Encoding}), com vocabulário de 50.257 \textit{tokens}. Manter o
tokenizador original é indispensável: os \textit{embeddings} do modelo foram aprendidos
para esse vocabulário específico, e qualquer substituição invalidaria os pesos
pré-treinados.

\subsection{Divisão dos Dados}\label{subsec:divisao_dados}

Adota-se a divisão oficial do corpus. O \textbf{conjunto de teste} é utilizado em todas as
medições de perplexity reportadas, tanto do modelo denso quanto de todas as configurações
comprimidas. O \textbf{conjunto de treinamento} é utilizado exclusivamente na etapa de
\textit{fine-tuning} de recuperação (Seção~\ref{sec:finetuning}). O \textbf{conjunto de
validação} é empregado apenas para acompanhar a evolução do treinamento ao longo dos
passos, quando essa curva é solicitada, e nunca para selecionar resultados reportados.

Essa separação estrita importa: como as varreduras \textit{one-shot} não realizam
treinamento algum, elas não têm contato com o conjunto de treinamento, e a perplexity de
teste é uma medida limpa do dano causado pela poda.

\section{Modelo de Referência}\label{sec:modelo}

\subsection{O GPT-2 \textit{small}}\label{subsec:gpt2small}

O modelo de referência é o GPT-2 \textit{small}, proposto por \citeonline{Radford2019} e
descrito na Seção~\ref{subsec:llm_gpt2}. Trata-se de um modelo autorregressivo composto
exclusivamente pelo componente decodificador da arquitetura \textit{Transformer}. Os pesos
pré-treinados são obtidos pela plataforma Hugging Face, que os disponibiliza publicamente,
e carregados por meio da biblioteca PyTorch \cite{Paszke2019}.

A Tabela~\ref{tab:gpt2_config} resume a configuração da variante utilizada.

\begin{table}[htb]
\centering
\caption{Configuração do GPT-2 \textit{small} utilizado como modelo de referência.}
\label{tab:gpt2_config}
\small
\begin{tabular}{p{6.5cm} r}
\hline
\textbf{Característica} & \textbf{Valor} \\
\hline
Blocos \textit{Transformer} empilhados & 12 \\
Cabeças de atenção por bloco & 12 \\
Dimensão do modelo ($d_{\text{model}}$) & 768 \\
Dimensão intermediária do MLP ($d_{\text{ff}}$) & 3.072 \\
Tamanho do vocabulário & 50.257 \\
Janela de contexto máxima & 1.024 \textit{tokens} \\
Total de parâmetros & 124.439.808 \\
Tamanho em memória & 474,70 MB \\
FLOPs por janela de 512 \textit{tokens} & 126,6 G \\
Perplexity no WikiText-2 (teste) & 25,17 \\
\hline
\end{tabular}\\
\textit{Fonte: Elaboração própria.}
\end{table}

\subsection{Justificativa da Escolha}\label{subsec:justificativa_gpt2}

O GPT-2 \textit{small} atende simultaneamente aos cinco critérios estabelecidos na
Seção~\ref{sec:metodologia_modelo}. É um modelo de linguagem representativo e amplamente
estudado; dispõe de pesos pré-treinados de acesso público; possui 124 milhões de
parâmetros, volume que cabe confortavelmente na memória da GPU disponível mesmo durante o
\textit{fine-tuning}; é frequentemente empregado como objeto de estudo em trabalhos de
compressão; e conta com implementação madura em PyTorch.

Cabe observar que, embora seja a menor variante da família GPT-2, o modelo é cerca de
cinco vezes maior que a ResNet-50, arquitetura tradicionalmente usada em estudos de poda
em visão computacional. Essa diferença de escala é justamente o que torna o trabalho
pertinente: as conclusões da literatura de poda foram estabelecidas majoritariamente em
redes convolucionais, e não é evidente que se transfiram para modelos com esse perfil de
parâmetros.

\subsection{Parâmetros Podáveis}\label{subsec:parametros_podaveis}

Uma decisão de implementação que condiciona todos os resultados é a definição de quais
parâmetros são elegíveis à poda. Adotam-se, por convenção consolidada na literatura,
apenas as matrizes de peso das projeções lineares internas aos blocos
\textit{Transformer}, a saber, quatro por bloco:

\begin{itemize}
    \item \texttt{attn.c\_attn} --- projeção que produz consulta, chave e valor;
    \item \texttt{attn.c\_proj} --- projeção de saída da atenção;
    \item \texttt{mlp.c\_fc} --- primeira transformação linear do MLP;
    \item \texttt{mlp.c\_proj} --- segunda transformação linear do MLP.
\end{itemize}

Com 12 blocos, são 48 matrizes podáveis. Ficam deliberadamente fora do escopo da poda:

\begin{itemize}
    \item as matrizes de \textit{embedding} de \textit{token} (\texttt{wte}) e de posição
          (\texttt{wpe}), que codificam o vocabulário e não representam computação
          repetida a cada camada;
    \item a cabeça de linguagem (\texttt{lm\_head}), que no GPT-2 compartilha pesos com
          \texttt{wte} --- podá-la corromperia os \textit{embeddings};
    \item vieses e parâmetros de normalização de camada, que são numericamente baratos e
          desproporcionalmente sensíveis.
\end{itemize}

Essa restrição tem consequência direta na leitura dos resultados: quando se afirma que uma
configuração tem 50\% de esparsidade, trata-se de 50\% das matrizes podáveis, não do
modelo inteiro. Como os \textit{embeddings} do GPT-2 concentram parcela expressiva dos
parâmetros, a esparsidade sobre o total é sempre menor que a esparsidade sobre o conjunto
podável, e as duas grandezas são reportadas separadamente.

\section{Métodos de Poda da Literatura}\label{sec:metodos_poda}

\subsection{Método 1 --- Poda Não Estruturada por Magnitude}\label{subsec:metodo_magnitude}

O primeiro método segue \citeonline{Han2015}. O critério de importância é o valor absoluto
do peso: assume-se que pesos próximos de zero contribuem pouco para a saída e podem ser
anulados. Dada uma taxa de esparsidade alvo, calcula-se o limiar correspondente ao
quantil apropriado da distribuição de $|w|$ e zeram-se todos os pesos abaixo dele.

A implementação preserva a característica essencial do método: os pesos são
\textbf{zerados}, não removidos. As matrizes conservam suas dimensões originais e a
arquitetura permanece intacta --- o que, como discutido na Seção~\ref{sec:pruning}, é
precisamente a razão pela qual esse método não reduz o custo de inferência em
\textit{hardware} convencional.

\subsection{Método 2 --- Poda Estruturada por Norma $L_1$}\label{subsec:metodo_estruturado}

O segundo método adapta \citeonline{Li2017} ao \textit{Transformer}. O trabalho original
poda filtros convolucionais ordenados por norma $L_1$; aqui, as unidades estruturais
equivalentes são as \textbf{cabeças de atenção} e os \textbf{neurônios da camada
intermediária do MLP}. Para cada uma dessas estruturas calcula-se a norma $L_1$ dos pesos
que lhe correspondem, e removem-se as de menor norma.

A diferença decisiva em relação ao método anterior é que a remoção é \textbf{física}. As
matrizes são efetivamente reconstruídas com dimensões menores: as colunas de consulta,
chave e valor associadas às cabeças removidas são eliminadas de \texttt{c\_attn}, as
linhas correspondentes saem de \texttt{c\_proj}, e a contabilidade interna do módulo de
atenção --- número de cabeças e dimensão de divisão --- é atualizada de acordo. O
resultado não é um modelo esparso, mas uma arquitetura menor e densa.

A implementação admite três alvos: apenas cabeças, apenas neurônios MLP, ou ambos
simultaneamente. Uma salvaguarda impede que qualquer bloco fique sem nenhuma cabeça ou
sem nenhum neurônio, situação que romperia o fluxo residual e inviabilizaria o modelo.

Convém registrar que a esparsidade alvo, neste método, é expressa como fração de
\textbf{estruturas} removidas, não de parâmetros. Duas razões fazem a fração de parâmetros
efetivamente eliminada divergir da taxa solicitada, e por isso ela é reportada
separadamente em todas as tabelas. A primeira é que cabeças e neurônios têm custos
diferentes em parâmetros. A segunda é a granularidade discreta das estruturas: o número de
unidades a remover é a parte inteira do produto entre a taxa e a contagem disponível, o que
nos neurônios do MLP --- 3.072 por bloco --- é praticamente imperceptível, mas nas cabeças
de atenção, que são apenas 12 por bloco, produz arredondamentos visíveis. Uma taxa de
90\%, por exemplo, remove 10 das 12 cabeças, ou seja, 83\%.

\subsection{Escopos de Alocação}\label{subsec:escopos}

Ambos os métodos são avaliados sob dois escopos, que representam formas distintas de
distribuir o orçamento de poda entre as camadas:

\begin{itemize}
    \item \textbf{Global:} o critério de importância é aplicado ao modelo inteiro de uma
          só vez. Um único limiar --- ou uma única ordenação de estruturas --- decide o
          que sai, de modo que camadas com pesos sistematicamente menores cedem
          proporcionalmente mais. Na poda estruturada com os dois alvos simultâneos, o
          escopo global compreende duas ordenações independentes, uma de cabeças e outra
          de neurônios, cada uma abrangendo as doze camadas: as normas $L_1$ de uma
          cabeça de atenção e de um neurônio do MLP não são grandezas comparáveis, pois
          agregam números distintos de pesos, e ordená-las em conjunto faria a escolha
          depender do tamanho da estrutura em vez de sua importância relativa.
    \item \textbf{Uniforme:} cada camada recebe exatamente a mesma fração de poda,
          decidida isoladamente dentro dela. É a alocação ingênua por excelência, pois
          ignora qualquer diferença de importância entre camadas.
\end{itemize}

O contraste entre esses dois escopos é o que estabelece a linha de base contra a qual a
estratégia otimizada é julgada. Adianta-se que nenhum deles domina o outro em toda a faixa
de esparsidade, e é essa constatação que motiva a análise de sensibilidade.

\subsection{Taxas de Poda Avaliadas}\label{subsec:taxas}

Cada combinação de método e escopo é avaliada em dez níveis: 0\% (o modelo denso, incluído
como controle interno de cada varredura) e de 10\% a 90\% em passos de 10 pontos
percentuais. A varredura completa da faixa é uma escolha deliberada: como não se conhecia
de antemão onde estaria o ponto de colapso de cada estratégia, restringir os níveis a uma
faixa estreita poderia ocultar exatamente o comportamento que se deseja caracterizar.

Cada nível parte de um modelo recarregado do zero, e não do modelo podado no nível
anterior. Sem esse cuidado, a poda seria cumulativa e os níveis deixariam de ser
independentes entre si.

\section{Estratégia Otimizada: Sensibilidade por Camada}\label{sec:estrategia_otimizada}

Os métodos descritos até aqui compartilham uma limitação: nenhum deles mede o quanto cada
camada realmente importa. O escopo uniforme trata todas as camadas como intercambiáveis; o
escopo global decide com base na magnitude dos pesos, que é um \textit{proxy} de
importância, não uma medida dela. A estratégia proposta substitui esse \textit{proxy} por
uma medição direta, em duas fases.

\subsection{Fase 1 --- Perfil de Sensibilidade}\label{subsec:fase1}

Para cada camada $\ell$ do modelo e para cada estratégia de poda, poda-se
\textbf{exclusivamente} aquela camada a uma taxa-sonda fixa, mantendo todas as demais
intactas, e mede-se a perplexity resultante. A degradação é registrada em escala
logarítmica:

\begin{equation}\label{eq:delta_perfil}
\Delta_{\ell} = \ln\!\left(\frac{\text{PPL}_{\ell}}{\text{PPL}_{\text{base}}}\right),
\end{equation}

\noindent onde $\text{PPL}_{\ell}$ é a perplexity do modelo com apenas a camada $\ell$
podada. O logaritmo é necessário porque a perplexity de modelos podados sem recuperação
varia por ordens de grandeza; em escala linear, uma única camada catastrófica tornaria
todas as demais indistinguíveis entre si.

Adota-se taxa-sonda de 50\%. O procedimento não envolve treinamento e produz, ao final, um
perfil com uma medida de sensibilidade por camada, para cada estratégia. Como as duas
estratégias removem coisas diferentes, os perfis são levantados separadamente e não são
intercambiáveis --- ponto que a Seção~\ref{subsec:perfis_divergem} demonstra
experimentalmente.

Como a escolha da taxa-sonda é arbitrária, o perfilamento completo é replicado a uma
segunda taxa --- 30\% ---, nas duas estratégias, de modo a verificar se a ordenação de
sensibilidade obtida é propriedade das camadas ou artefato do regime em que a sonda opera.
A réplica cumpre função de verificação e não alimenta nenhuma alocação: todas as taxas por
camada reportadas derivam do perfil de 50\%. O resultado da verificação é apresentado na
Seção~\ref{subsec:robustez_sonda}.

\subsection{Fase 2 --- Alocação Adaptativa de Taxas}\label{subsec:fase2}

O perfil é então convertido em uma taxa de poda por camada. Atribui-se a cada camada uma
taxa proporcional à sua robustez --- o inverso da degradação medida --- elevada a um
expoente $\beta$ e escalada por um fator único $\alpha$, escolhido de modo que a média das
taxas, ponderada pelo número de parâmetros podáveis de cada camada, iguale o orçamento
global $s$ pretendido:

\begin{equation}\label{eq:alocacao_dev}
r_{\ell} = \min\!\left(\alpha \cdot \Delta_{\ell}^{-\beta},\ r_{\max}\right),
\qquad
\frac{\sum_{\ell} w_{\ell}\, r_{\ell}}{\sum_{\ell} w_{\ell}} = s,
\end{equation}

\noindent em que $w_{\ell}$ é o número de parâmetros podáveis da camada $\ell$ e
$r_{\max} = 0{,}95$ é um teto que impede que mesmo a camada mais robusta seja
integralmente zerada. Como a fração podada é monótona em $\alpha$, o fator é obtido por
bisseção.

A ponderação por $w_{\ell}$ não produz efeito no modelo aqui empregado: os doze blocos do
GPT-2 são idênticos em estrutura e, portanto, têm exatamente o mesmo número de parâmetros
podáveis, de modo que a média ponderada se reduz à média simples das taxas. Ela é mantida
na formulação por generalidade --- em arquiteturas de blocos heterogêneos, ignorá-la faria
o orçamento global ser sistematicamente descumprido.

O expoente $\beta$ controla o quanto o perfil é levado a sério. Com $\beta = 1$ emprega-se
o inverso puro da degradação; valores intermediários amortecem a razão entre camadas. Esse
amortecimento não é um detalhe: a razão de robustez entre duas camadas é ilimitada, e sem
controle uma única camada muito robusta absorveria praticamente todo o orçamento. Os
experimentos avaliam $\beta = 1{,}0$ e $\beta = 0{,}5$ sobre o mesmo perfil, o que fornece
a ablação da função de alocação sem exigir a repetição da fase 1, que é a etapa cara.

O valor $\beta = 0$ é avaliado também, mas com outra finalidade: como o fator
$\Delta_{\ell}^{-\beta}$ vale 1 para toda camada, todas recebem a mesma taxa e a regra deve
reproduzir a poda uniforme. Serve, portanto, de controle de corretude da alocação, e não de
alocação candidata. O controle é exato na poda estruturada; na poda por magnitude ele não
se aplica, porque a poda por camada adota um limiar único por bloco --- compartilhado pelas
quatro projeções --- enquanto a varredura uniforme decide um limiar por matriz. As duas
alcançam a mesma esparsidade agregada, distribuída de forma distinta dentro do bloco. A
verificação é reportada na Seção~\ref{subsec:controle_beta_zero}.

Registre-se, por fim, uma salvaguarda da implementação não visível na
Equação~\ref{eq:alocacao_dev}: como a degradação medida pela sonda pode ser nula ou
negativa --- há casos em que remover cabeças de atenção \textit{melhora} a perplexity,
fenômeno reportado por \citeonline{Michel2019} ---, a robustez $\Delta_{\ell}^{-1}$ é
calculada sobre um piso positivo, de modo que tais camadas entrem na alocação como as mais
robustas em vez de produzir divisão por zero. Nos perfis efetivamente levantados neste
trabalho o piso não foi acionado: a menor degradação observada é de 0,0346, duas ordens de
grandeza acima dele.

\section{\textit{Fine-tuning} de Recuperação}\label{sec:finetuning}

Todas as etapas anteriores operam em regime \textit{one-shot}. Esse regime isola o dano
causado pela poda, mas não corresponde ao procedimento da literatura: tanto
\citeonline{Han2015} quanto \citeonline{Li2017} retreinam o modelo após podá-lo. Esta
etapa fecha essa lacuna, medindo, para cada configuração, a perplexity antes e depois do
retreinamento.

Duas decisões de implementação merecem registro por afetarem a validade dos resultados.

A primeira é o \textbf{controle denso}. O GPT-2 não foi pré-treinado no WikiText-2, de
modo que o simples \textit{fine-tuning} do modelo denso nesse corpus já reduz a perplexity
por adaptação de domínio, sem poda alguma envolvida. Sem essa linha de controle, seria
impossível separar a recuperação do dano da poda da mera adaptação ao corpus. O modelo
denso é, portanto, retreinado com o mesmo orçamento de passos e avaliado junto às demais
configurações.

A segunda é a \textbf{preservação da esparsidade}. Na poda não estruturada os pesos são
apenas zerados, e o gradiente os restauraria já no primeiro passo do otimizador; ao final
do treinamento, o modelo estaria novamente denso. Como em \citeonline{Han2015}, mantém-se
uma máscara binária fixa, capturada logo após a poda e reaplicada depois de cada passo do
otimizador. Registre-se que zerar o gradiente não seria suficiente: o otimizador utilizado
possui momento e regularização por decaimento de pesos, e ambos alteram um parâmetro mesmo
quando seu gradiente é nulo. A esparsidade medida ao final do treinamento é registrada
como verificação de integridade. Na poda estruturada não há máscara, pois as estruturas já
foram fisicamente removidas.

O treinamento emprega o otimizador AdamW, com aquecimento linear da taxa de aprendizado
seguido de decaimento linear, recorte de norma do gradiente e acúmulo de gradiente para
viabilizar lotes efetivos maiores do que a memória da GPU comportaria. O decaimento de
pesos é aplicado apenas às matrizes de peso, deixando de fora vieses e parâmetros de
normalização, conforme prática consolidada. A passagem direta e a retropropagação executam
em \textbf{precisão mista}, com escalonamento dinâmico da perda para evitar
subtransbordamento dos gradientes em meia precisão; os pesos são mantidos em precisão
simples. Essa escolha reduz aproximadamente à metade a memória de ativações e é o que
viabiliza o orçamento de treino na GPU disponível, mas introduz uma fonte adicional de
variação numérica que não existe nas etapas \textit{one-shot}, todas executadas em
precisão simples. O consumo de
energia do treinamento é medido separadamente do consumo da inferência: trata-se de um
custo pago uma única vez, que precisa ser amortizado contra a economia por inferência para
que o balanço energético do trabalho feche corretamente.

Diferentemente de todas as etapas anteriores, esta é estocástica --- a ordem dos lotes e o
\textit{dropout} introduzem variabilidade --- e, portanto, é a única em que a semente
aleatória efetivamente altera o resultado.

\section{\textit{Framework} de Implementação}\label{sec:framework}

A implementação é feita em Python, com o \textit{framework} PyTorch \cite{Paszke2019}
para construção, manipulação e treinamento dos modelos. O PyTorch foi escolhido por seu
sistema de diferenciação automática, que constrói dinamicamente o grafo computacional e
simplifica a implementação do \textit{backpropagation} durante o \textit{fine-tuning}, e
por seu suporte nativo à aceleração em GPU.

Os pesos pré-treinados são obtidos pela plataforma Hugging Face por meio da biblioteca
\texttt{transformers}. Cabe a distinção, já feita no Capítulo~\ref{chap:metodologia}: a
Hugging Face é a plataforma que hospeda e distribui os modelos, ao passo que PyTorch é a
biblioteca que os executa.

As demais bibliotecas empregadas são:

\begin{itemize}
    \item \texttt{datasets}: carregamento do corpus WikiText a partir da plataforma
          Hugging Face;
    \item \texttt{thop}: contagem de operações de ponto flutuante por análise do grafo
          computacional;
    \item \texttt{CodeCarbon} \cite{CodeCarbon2023}: medição empírica do consumo de
          energia e estimativa das emissões de $CO_2$, a partir do monitoramento do
          \textit{hardware} combinado à intensidade de carbono da rede elétrica local.
\end{itemize}

Uma observação relevante sobre a implementação da poda: ambas as estratégias foram
implementadas manualmente, sem recurso a bibliotecas dedicadas. No caso da poda
estruturada, essa decisão foi imposta pelo ambiente --- as versões mais recentes da
biblioteca \texttt{transformers} removeram a interface de poda de cabeças que existia em
versões anteriores. A remoção física das estruturas é, portanto, realizada diretamente:
as matrizes são reconstruídas com as dimensões reduzidas e a contabilidade interna do
módulo de atenção é atualizada. A implementação da poda é coberta por uma
suíte de testes automatizados (\texttt{test\_poda}), executada sobre um GPT-2 de dimensões
reduzidas e pesos aleatórios --- o que está sob verificação é a mecânica da poda, não os
pesos pré-treinados. O teste central é o de \textbf{equivalência numérica}: remover
fisicamente um conjunto de cabeças de atenção, ou de neurônios do MLP, deve produzir
exatamente as mesmas saídas que apenas anular os pesos correspondentes no modelo de
dimensões originais. A equivalência é matematicamente necessária --- anular as fatias de
consulta, chave e valor de uma cabeça zera seu vetor de valores, de modo que sua
contribuição à saída da atenção se anula qualquer que seja a distribuição de atenção
calculada ---, e por isso qualquer erro de indexação na reconstrução das matrizes a
quebraria. Os demais testes verificam que a contabilidade interna do módulo de atenção
permanece coerente com as formas dos tensores, que a salvaguarda preserva ao menos uma
estrutura por bloco mesmo sob taxa integral, que as estruturas removidas são efetivamente
as de menor norma $L_1$, que a poda por magnitude atinge a esparsidade solicitada
preservando as dimensões, que a regra de alocação respeita o orçamento e poda menos a
camada mais sensível, e que a máscara de esparsidade sobrevive a um passo do otimizador.

\section{\textit{Pipeline} Experimental}\label{sec:pipeline}

O código está organizado em módulos de responsabilidade única, o que permite que cada
estratégia de poda reutilize exatamente a mesma rotina de avaliação --- condição para que
as comparações sejam legítimas. A Tabela~\ref{tab:modulos} descreve os módulos e a
Figura~\ref{fig:pipeline_experimental} representa o fluxo experimental.

\begin{table}[htb]
\centering
\caption{Organização modular da implementação.}
\label{tab:modulos}
\small
\begin{tabular}{p{3.4cm} p{9.6cm}}
\hline
\textbf{Módulo} & \textbf{Responsabilidade} \\
\hline
\texttt{config} & Configuração central: modelo, corpus, parâmetros de avaliação, semente e ambiente. \\
\texttt{utils} & Fixação das sementes e escrita incremental dos resultados em arquivo. \\
\texttt{data} & Carregamento e tokenização do corpus. \\
\texttt{model} & Carregamento do GPT-2 pré-treinado. \\
\texttt{eval} & Perplexity por janela deslizante e métricas de custo computacional. \\
\texttt{energy} & Medição de energia e emissões. \\
\texttt{prune\_magnitude} & Poda não estruturada por magnitude, nos dois escopos. \\
\texttt{prune\_structured} & Poda estruturada de cabeças e neurônios, com remoção física. \\
\texttt{sensitivity} & Perfil de sensibilidade (fase 1) e alocação adaptativa (fase 2). \\
\texttt{finetune} & \textit{Fine-tuning} de recuperação, com preservação de esparsidade. \\
\texttt{run} & Orquestração da avaliação do modelo de referência. \\
\texttt{test\_poda} & Testes de corretude da poda, descritos na Seção~\ref{sec:framework}. \\
\hline
\end{tabular}\\
\textit{Fonte: Elaboração própria.}
\end{table}

\begin{figure}[htb]
\centering
\caption{Fluxo experimental do trabalho. Todas as estratégias convergem para a mesma
rotina de avaliação, o que garante a simetria necessária às comparações.}
\label{fig:pipeline_experimental}
\begin{tikzpicture}[font=\scriptsize, node distance=0.55cm and 0.45cm,
  fonte/.style={rectangle,rounded corners=2pt,draw,thick,fill=gray!8,
                minimum width=3.1cm,minimum height=0.75cm,align=center},
  etapa/.style={rectangle,rounded corners=2pt,draw,thick,fill=gray!18,
                minimum width=3.3cm,minimum height=0.9cm,align=center},
  saida/.style={rectangle,rounded corners=2pt,draw,thick,fill=gray!30,
                minimum width=4.6cm,minimum height=0.75cm,align=center}]

  % --- entradas ---
  \node[fonte] (mod) at (-2.1,0) {GPT-2 \textit{small}\\ \scriptsize pré-treinado};
  \node[fonte] (dat) at (2.1,0) {WikiText-2\\ \scriptsize corpus de avaliação};

  % --- baseline ---
  \node[etapa] (base) at (0,-1.35) {\textbf{Modelo de referência}\\ \scriptsize perplexity 25,17};
  \draw[-{Stealth}] (mod.south) -- ++(0,-0.3) -| (base.north);
  \draw[-{Stealth}] (dat.south) -- ++(0,-0.3) -| (base.north);

  % --- três estratégias ---
  \node[etapa] (mag) at (-4.6,-3.0) {\textbf{Magnitude}\\ \scriptsize global / uniforme};
  \node[etapa] (est) at (0,-3.0)    {\textbf{Estruturada}\\ \scriptsize global / uniforme};
  \node[etapa] (sen) at (4.6,-3.0)  {\textbf{Sensibilidade}\\ \scriptsize perfil $\rightarrow$ alocação};
  \foreach \n in {mag,est,sen} \draw[-{Stealth}] (base.south) -- ++(0,-0.35) -| (\n.north);

  % --- fine-tuning ---
  \node[etapa] (ft) at (0,-4.75) {\textbf{\textit{Fine-tuning} de recuperação}\\
        \scriptsize com controle denso e máscara de esparsidade};
  \foreach \n in {mag,est,sen} \draw[-{Stealth}] (\n.south) -- ++(0,-0.35) -| (ft.north);

  % --- avaliação ---
  \node[etapa] (av) at (0,-6.3) {\textbf{Avaliação comum}\\
        \scriptsize perplexity, parâmetros, tamanho, FLOPs, tempo, energia};
  \draw[-{Stealth}] (ft) -- (av);

  \node[saida] (out) at (0,-7.6) {Tabelas e figuras do Capítulo~\ref{chap:resultados}};
  \draw[-{Stealth}] (av) -- (out);

  % --- desvio one-shot ---
  \draw[-{Stealth},dashed,gray!70] (mag.west) -- ++(-0.55,0)
     |- node[left=2pt,align=right,font=\scriptsize,black,pos=0.25]
        {regime\\ \textit{one-shot}} (av.west);
\end{tikzpicture}
\fonte{Elaborada pelo autor.}
\end{figure}

Como indica a linha tracejada da figura, as varreduras \textit{one-shot} contornam a etapa
de \textit{fine-tuning} e vão direto à avaliação. Essa bifurcação é o que permite separar,
nos resultados, o dano causado pela poda daquilo que o retreinamento consegue recuperar.

\section{Métricas de Avaliação}\label{sec:metricas}

Toda configuração --- o modelo denso, as podadas pelos métodos da literatura e as podadas
pela estratégia otimizada --- é avaliada pelo mesmo conjunto de métricas, cobrindo as três
dimensões definidas na Seção~\ref{sec:metodologia_metricas}. A
Tabela~\ref{tab:metricas} as resume.

\begin{table}[htb]
\centering
\caption{Métricas de avaliação coletadas para cada configuração.}
\label{tab:metricas}
\small
\begin{tabular}{p{3.2cm} p{6.1cm} p{1.7cm} p{2.2cm}}
\hline
\textbf{Métrica} & \textbf{O que mede} & \textbf{Unidade} & \textbf{Ferramenta} \\
\hline
Perplexity & Qualidade da modelagem da linguagem; quanto menor, melhor & --- & PyTorch \\
Parâmetros totais & Parâmetros existentes na arquitetura & unidades & PyTorch \\
Parâmetros não nulos & Parâmetros efetivamente diferentes de zero & unidades & PyTorch \\
Esparsidade real & Fração de parâmetros podáveis eliminados & \% & PyTorch \\
Tamanho do modelo & Memória ocupada pelos parâmetros & MB & PyTorch \\
FLOPs & Operações de ponto flutuante por inferência & FLOPs & \texttt{thop} \\
Tempo de inferência & Tempo médio de uma passagem direta & ms & \texttt{perf\_counter} \\
Energia consumida & Energia gasta durante a avaliação & kWh & \texttt{CodeCarbon} \\
Emissões de $CO_2$ & Emissões equivalentes ao consumo & kg & \texttt{CodeCarbon} \\
\hline
\end{tabular}\\
\textit{Fonte: Elaboração própria.}
\end{table}

\subsection{Perplexity}\label{subsec:metrica_perplexity}

A perplexity é a métrica de desempenho preditivo adotada. Ela é calculada pelo método da
janela deslizante, padrão para modelos autorregressivos: o fluxo de \textit{tokens} do
conjunto de teste é percorrido em janelas do tamanho máximo de contexto do modelo, com
passo menor que a janela, de modo que cada \textit{token} seja previsto dispondo de
contexto anterior substancial. Apenas os \textit{tokens} ainda não avaliados contribuem
para a perda em cada janela; os demais entram como contexto.

Emprega-se janela de 1.024 \textit{tokens} com passo de 512. Passos menores produzem
estimativas mais precisas ao custo de mais computação; o valor adotado é o compromisso
usual na literatura, e é mantido idêntico em todas as configurações avaliadas.

\subsection{Custo Computacional}\label{subsec:metrica_custo}

A contagem de \textbf{parâmetros} é reportada em duas formas complementares. O total
reflete o que existe na arquitetura; a contagem de não nulos reflete o que sobrou depois da
poda. Na poda não estruturada as duas divergem --- há zeros, mas a arquitetura é a mesma
---, enquanto na poda estruturada elas coincidem, pois o que foi removido deixou de
existir. Essa divergência é, por si só, um resultado: ela evidencia a diferença entre
compressão contábil e compressão física.

A \textbf{esparsidade real} reportada nas tabelas é sempre calculada sobre o conjunto
podável, mas a natureza do cálculo acompanha a estratégia: na poda por magnitude é a fração
de pesos anulados nas 48 matrizes; na poda estruturada é a fração de parâmetros que
deixaram de existir nas mesmas projeções, vieses incluídos, já que a remoção de uma
estrutura elimina também o viés que lhe corresponde. As duas grandezas medem a mesma coisa
--- quanto do conjunto podável foi eliminado ---, mas não são idênticas em denominador, e
diferenças na segunda casa decimal entre estratégias não devem ser interpretadas.

Os \textbf{FLOPs} são contabilizados por análise do grafo computacional. Registre-se uma
particularidade da implementação: as projeções internas do GPT-2 não são camadas lineares
convencionais, e sim um tipo específico de camada que a ferramenta de contagem não
reconhece por padrão. Sem tratamento, a contagem ignoraria justamente as matrizes que a
poda modifica, e a métrica sairia constante. Foi necessário, portanto, registrar uma regra
de contagem própria para esse tipo de camada. Reconhece-se uma limitação remanescente: as
multiplicações de matriz internas ao mecanismo de atenção não são contabilizadas pela
ferramenta em nenhuma configuração. Como a omissão é idêntica para todas as estratégias,
as comparações permanecem válidas. A implementação prevê ainda um estimador analítico de
reserva --- duas operações por parâmetro não pertencente aos \textit{embeddings}, por
\textit{token} processado ---, acionado apenas se a contagem pelo grafo falhar. Todos os
valores reportados no Capítulo~\ref{chap:resultados} provêm da contagem pelo grafo, o que
se verifica pelo fato de os FLOPs variarem com a poda estruturada; o estimador de reserva
produziria uma grandeza de natureza diferente e não deve ser misturado a eles.

O \textbf{tempo de inferência} é medido sobre uma entrada representativa de 512
\textit{tokens}, com execuções de aquecimento descartadas --- para mitigar efeitos de
inicialização da GPU --- seguidas de execuções cronometradas, das quais se reporta média e
desvio-padrão.

\subsection{Impacto Ambiental}\label{subsec:metrica_energia}

O consumo de \textbf{energia} é medido com o \texttt{CodeCarbon}, que monitora o
\textit{hardware} durante a execução do bloco de avaliação e reporta a energia consumida em
quilowatts-hora. As \textbf{emissões de $CO_2$} são estimadas a partir desse consumo e da
intensidade de carbono da rede elétrica onde o código executa.

Sobre esta última métrica cabe uma ressalva importante, identificada durante os
experimentos: a intensidade de carbono utilizada varia conforme a sessão de execução. Em
consequência, as emissões não são comparáveis entre execuções realizadas em sessões
distintas, e a análise dos resultados apoia-se no consumo de energia, não nas emissões
derivadas dele.

\section{Ambiente de Execução e Reprodutibilidade}\label{sec:ambiente}

Os experimentos são executados na plataforma Kaggle, em sessões com GPU NVIDIA Tesla T4. A
implementação detecta automaticamente o dispositivo disponível, o que permite executar
localmente em CPU para testes rápidos sem alteração de código.

Quanto à reprodutibilidade, três aspectos merecem registro.

O primeiro é que as varreduras \textit{one-shot} são \textbf{determinísticas}. Não há
treinamento, o modelo opera em modo de avaliação, o conjunto de teste é fixo e os critérios
de poda são ordenações determinísticas sobre pesos fixos. Verificou-se empiricamente que
execuções independentes, em sessões distintas, reproduzem a perplexity até a última casa
decimal. Repetições com sementes diferentes, nessas etapas, produziriam linhas idênticas e
não acrescentariam informação. A semente passa a importar apenas no \textit{fine-tuning}.

O segundo é que as medições de \textbf{tempo e energia dependem da sessão}. Observou-se que
o mesmo modelo denso, sem qualquer poda, registrou tempos e consumos sensivelmente
distintos em sessões diferentes da plataforma, o que se atribui à variação de
\textit{hardware} físico e de carga concorrente. Por essa razão, comparações de custo entre
varreduras executadas em sessões distintas são sempre feitas em termos relativos ao modelo
denso da própria sessão, nunca em valores absolutos.

O terceiro é que cada configuração experimental grava uma linha em arquivo de resultados,
contendo não apenas as métricas mas também os parâmetros que a geraram --- estratégia,
escopo, taxa alvo, taxa efetiva e, no caso da estratégia otimizada, as taxas atribuídas a
cada camada. Isso permite rastrear qualquer número apresentado no
Capítulo~\ref{chap:resultados} até a configuração exata que o produziu.

\section{Limitações Metodológicas}\label{sec:limitacoes}

Reconhecem-se, desde o desenho experimental, limitações que condicionam o escopo dos
resultados:

\begin{itemize}
    \item \textbf{Um único modelo e um único corpus.} Os experimentos empregam
          exclusivamente o GPT-2 \textit{small} sobre o WikiText-2. Os resultados podem
          não se transferir diretamente para modelos de escala substancialmente maior ou
          para corpora de natureza distinta. A implementação foi mantida parametrizável
          justamente para viabilizar essa verificação em trabalhos futuros.

    \item \textbf{Duas estratégias de poda, e não mais.} A restrição de tempo experimental
          levou à escolha de duas estratégias que cobrem os extremos de granularidade.
          Métodos baseados em informação de gradiente ou de segunda ordem ficaram fora do
          escopo.

    \item \textbf{Faixa restrita de taxas-sonda no perfil de sensibilidade.} O perfil que
          alimenta a alocação é levantado a uma única taxa, e sua estabilidade foi
          verificada por meio de uma réplica a uma segunda taxa
          (Seção~\ref{subsec:robustez_sonda}). Ambas pertencem, porém, ao mesmo regime
          intermediário de poda: nada garante que a ordenação de sensibilidade entre
          camadas se mantenha em regimes muito mais brandos ou muito mais agressivos.

    \item \textbf{Sensibilidade medida em isolamento.} A sonda mede o efeito de cada camada
          isoladamente e ignora interações entre elas. Como se verifica no
          Capítulo~\ref{chap:resultados}, a degradação é fortemente superaditiva, de modo
          que o perfil deve ser lido como ordenação relativa de importância, e não como
          modelo preditivo do desempenho de uma configuração completa.

    \item \textbf{Dependência do \textit{hardware}.} Tempo e energia são intrinsecamente
          dependentes da máquina. Os valores absolutos não devem ser comparados entre
          ambientes distintos; a análise apoia-se em comparações relativas dentro de um
          mesmo ambiente.

    \item \textbf{Hiperparâmetros não otimizados exaustivamente.} Os hiperparâmetros do
          \textit{fine-tuning} foram definidos a partir da literatura e de experimentos
          preliminares, sem busca sistemática. Valores melhores podem existir.

    \item \textbf{Semente única no \textit{fine-tuning}.} Esta é a única etapa estocástica
          do experimento, e todas as configurações foram executadas com uma única semente.
          As perplexities pós-retreinamento não vêm, portanto, acompanhadas de estimativa
          de variância, e diferenças pequenas entre configurações retreinadas devem ser
          lidas com a cautela correspondente. A comparação principal da etapa --- alocação
          guiada contra alocação global --- é afetada por essa limitação justamente no
          orçamento em que as duas praticamente empatam.

    \item \textbf{Estimativa de emissões.} A conversão de energia em emissões depende de
          tabelas regionais de intensidade de carbono, que carregam incerteza e variaram
          entre as sessões de execução. As emissões devem ser lidas como aproximações de
          ordem de grandeza.
\end{itemize}

Apesar dessas limitações, o desenho descrito é suficiente para responder às duas perguntas
enunciadas na Seção~\ref{sec:visao_geral}, que são o objeto do
Capítulo~\ref{chap:resultados}.
